% from ega0-19-fr.tex
\oldpage[0\textsubscript{IV-1}]{100}

  \item[\rm(iii)] Réciproquement, supposons que $B$ soit préadmissible, que la
  suite $(\mathfrak J^n)$ tende vers $0$, et que, pour tout $\lambda$,
  $(\mathfrak J/\mathfrak J^2)\otimes_C C_\lambda$ soit un $C_\lambda$-module
  projectif et l'homomorphisme (19.5.4.1) soit bijectif. Alors $B$ est une
  $A$-algèbre formellement lisse.
\end{enumerate}
\end{corollary}

\begin{proof}
En effet, la topologie de $\mathfrak J/\mathfrak J^2$ et celle des
$\mathbf S_C^n(\mathfrak J/\mathfrak J^2)$ sont alors aussi déduites de celle
de $B$ (19.5.1)~; les conclusions sont alors des conséquences immédiates de
(19.5.3) et des critères de (19.1.4) et (19.2.4) spéciaux à ce type de
topologies.
\end{proof}

\begin{remark}[19.5.8]
\label{0.19.5.8-fr}
Supposons que $\mathfrak J/\mathfrak J^2$ soit un $C$-module de type fini et
que, pour la topologie quotient de celle de $B$, $C$ soit un anneau de Zariski~;
soit $\mathfrak r$ un idéal de définition de $C$, de sorte que les
$\mathfrak r^n$ forment un système fondamental de voisinages de $0$ dans $C$.
Alors les conclusions de (i) et (ii) dans (19.5.7) peuvent être remplacées par
les suivantes~:
\begin{enumerate}
  \item[\rm(i')] $\mathfrak J/\mathfrak J^2$ est un $C$-module projectif.
  \item[\rm(ii')] L'homomorphisme canonique
  $\varphi:\mathbf S_C^*(\mathfrak J/\mathfrak J^2)
  \to\gr_{\mathfrak J}^*(B)$ est bijectif.
\end{enumerate}
En effet, il est clair que (i') entraîne la conclusion de (i) dans (19.5.7).
Inversement, si $(\mathfrak J/\mathfrak J^2)\otimes_C C_\lambda$ est un
$C_\lambda$-module projectif pour tout $\lambda$, alors
$(\mathfrak J/\mathfrak J^2)\otimes_C(C/\mathfrak r^n)$ est un
$(C/\mathfrak r^n)$-module projectif (donc plat) pour tout $n$~; on en conclut
que $\mathfrak J/\mathfrak J^2$ est un $C$-module plat
($\mathbf 0_{\mathrm{III}}$, 10.2.2), donc projectif puisqu'il est de
présentation finie (Bourbaki, \emph{Alg. comm.}, chap.~II, \textsection5,
n\textsuperscript{o}~2, cor.~2 du th.~1). D'autre part, les $C$-modules
$\mathbf S_C^n(\mathfrak J/\mathfrak J^2)$ et $\gr_{\mathfrak J}^n(B)$ sont
de type fini, et on sait que lorsque $C$ est un anneau de Zariski, il revient
au même de dire alors que $\varphi_n$ est bijectif ou que
$\widehat\varphi_n$ est bijectif (Bourbaki, \emph{Alg. comm.}, chap.~III,
\textsection3, n\textsuperscript{o}~5, prop.~9), donc (ii) est équivalent à
(ii').
\end{remark}

\subsection{Cas des algèbres sur un corps}
\label{subsection:0.19.6-fr}

\begin{theorem}[19.6.1) (Cohen]
\label{0.19.6.1-fr}
Soient $k$ un corps, $K$ une extension de $k$, $k$ et $K$
étant munis des topologies discrètes. Pour que $K$ soit une $k$-algèbre
formellement lisse, il faut et il suffit que $K$ soit une extension séparable
de $k$.
\end{theorem}

\begin{proof}
La nécessité de la condition sera établie en (19.6.5.1) (et naturellement ne
sera pas utilisée jusque-là)~; nous nous bornerons ici à prouver que la
condition est \emph{suffisante}. Distinguons deux cas~:

\emph{I. --- $K$ est une extension séparable de type fini de $k$.} On sait
alors (Bourbaki, \emph{Alg.}, chap.~V, \textsection9,
n\textsuperscript{o}~3, th.~2) qu'il existe une sous-extension pure
$K'=k(T_1,\ldots,T_n)$ de $K$ telle que $K$ soit une extension algébrique
finie séparable de $K'$. Compte tenu de (19.3.5, (iii)), on peut donc se
borner, soit au cas où $K=K'$, soit au cas où $K$ est algébrique fini sur $k$.
Dans le premier cas, on sait que $A=k[T_1,\ldots,T_n]$ est une $k$-algèbre
formellement lisse (19.3.3), donc il en est de même de
$k(T_1,\ldots,T_n)=S^{-1}A$, où $S=A-\{0\}$ (19.3.5, (iv)). Dans le second
cas, \emph{tous} les groupes de cohomologie de Hochschild
$\mathrm H_k^n(K,L)$, pour un $(K,K)$-bimodule quelconque $L$, sont nuls~: en
effet, si l'on considère la $k$-algèbre produit tensoriel $C=K\otimes_k K$,
$L$ est un $C$-module à gauche et le groupe de cohomologie
$\mathrm H_k^n(K,L)$ est égal à $\operatorname{Ext}_C^n(K,L)$, où $K$ est
aussi considéré comme $(K,K)$-bimodule,

